\documentclass[11pt]{article}

% ---- arXiv-style preprint preamble (standard packages only) ----
\usepackage[margin=1in]{geometry}
\usepackage{amsmath,amssymb}
\usepackage{booktabs}
\usepackage{array}
\usepackage{microtype}
\usepackage{xcolor}
\usepackage{siunitx}
\usepackage[hidelinks]{hyperref}
\usepackage{titlesec}
\usepackage{authblk}
\usepackage{caption}
\captionsetup{font=small,labelfont=bf}

\definecolor{accent}{HTML}{0D6B74}
\hypersetup{colorlinks=true,linkcolor=accent,citecolor=accent,urlcolor=accent}

\titleformat{\section}{\normalfont\large\bfseries\color{accent}}{\thesection}{0.6em}{}
\titleformat{\subsection}{\normalfont\normalsize\bfseries}{\thesubsection}{0.6em}{}

\newcommand{\dGhyd}{\ensuremath{\Delta G_{\mathrm{hyd}}}}
\newcommand{\Kwa}{\ensuremath{K_{\mathrm{wa}}}}
\newcommand{\kcal}{\,\mathrm{kcal\,mol^{-1}}}
\sisetup{detect-weight=true}

\title{\bfseries Reconciling a large, heterogeneous hydration free-energy dump:\\
a proportional-similarity consensus gate versus hierarchical partial pooling}

\author[1]{Guillaume Godin}
\affil[1]{Osmo Labs \quad\texttt{guillaume@osmo.ai}}
\date{\today}

\begin{document}
\maketitle

\begin{abstract}
\noindent
The GuthrieSolv database is the largest literature compilation of experimental
small-molecule hydration free energies (\num{53895} records, \num{5309} molecules),
but it is heterogeneous and of uncertain quality: raw measurements are stored in
172 distinct unit strings spanning vapour pressures, Henry's-law constants,
aqueous solubilities and free energies, and it ships with no data-reconciliation
code. We (i) harmonize the convertible fraction of the database into Ben--Naim
hydration free energies (\num{12027} observations over \num{2473} molecules,
temperature-corrected, keyed by InChIKey), validating every conversion route
against the curated FreeSolv benchmark; and (ii) use the resulting redundant,
multi-source, per-molecule observations to compare two families of reconciliation
method. The first is a \emph{proportional-similarity} (Homoset) gate that aligns
whole measurement sources by a median offset and accepts them into a consensus
only when a $\chi^2$-based concordance statistic exceeds a critical value; the
second is a \emph{mixle}-style hierarchical Gaussian partial-pooling model that
estimates a per-source bias, a per-source noise scale, and a shrunken per-molecule
posterior jointly by expectation--maximization, with an optional Student-$t$
robustification. On the 556-molecule FreeSolv overlap both methods substantially
outperform naïve pooling. The Homoset gate is the strongest transparent
\emph{curator}: when its noise scale $L$ is anchored to the true experimental
uncertainty it acts as a hard source-quality filter and attains a median-based
MAE of \SI{0.283}{\kcal}. The hierarchical model is the strongest
\emph{reconciler}: it attains the lowest RMSE (\SI{0.866}{\kcal}) and bias
(\SI{+0.07}{\kcal}), keeps all data, and uniquely recovers molecules that have no
clean reference measurement. We argue the two are complementary --- reconcile with
the hierarchical model, triage with the Homoset gate --- and release the
harmonized dataset and code.
\end{abstract}

\section{Introduction}
Experimental hydration (solvation-into-water) free energies underpin the
parameterization and validation of implicit-solvent models, physics-based free-energy
calculations, and machine-learned property predictors. The community gold standard,
FreeSolv~\cite{freesolv}, is small (\num{642} molecules) precisely because careful
curation is expensive. J.~P.~Guthrie assembled a far larger compilation from the primary
literature; after his death the family released it as GuthrieSolv~\cite{guthriesolv} with
explicit disclaimers about quality and an appeal for community curation. The database
mixes physical quantities and units within a single column, uses at least two hydration
standard states, and encodes data-trust flags in discrete uncertainty values. Crucially it
provides \emph{no} automated reconciliation --- its only quality-assurance code checks that
SMILES strings parse.

This is a natural test bed for multi-source data curation. Because many molecules appear in
GuthrieSolv multiple times --- across publications, temperatures and measurement techniques
--- each molecule carries a small, noisy, redundant set of observations once the raw values
are placed on a common scale. The curation question is then: given several disagreeing
measurements of the same quantity, what single value should one trust, and which molecules
should a human re-examine? We compare two answers with very different philosophies.

\section{Data harmonization}
\label{sec:harm}
GuthrieSolv stores a raw value (\texttt{value1}), a unit (\texttt{dimension1}), a temperature
(\texttt{te53p}) and, importantly, a measurement-type code (\texttt{process}) per record. The
process code is the key that disambiguates the 172 unit strings; Table~\ref{tab:process}
summarizes the convertible routes. We target the Ben--Naim standard state (transfer of one
mole per litre of ideal gas to one mole per litre of dilute aqueous solution), which is the
convention FreeSolv reports, so that converted values are directly comparable:
\begin{equation}
\dGhyd = -RT \ln \Kwa, \qquad \Kwa = C_{\mathrm{aq}}/C_{\mathrm{gas}} .
\label{eq:bennaim}
\end{equation}

\begin{table}[t]
\centering
\small
\begin{tabular}{@{}llrl@{}}
\toprule
\texttt{process} & measurement type & records & route to \dGhyd \\
\midrule
FEOH, DGS & free energy of hydration / solvation & \num{3160} & unit conversion (reference) \\
KWG, KGW & Henry's-law constant & \num{11751} & Eq.~\eqref{eq:bennaim} via \Kwa \\
VP $+$ AQSOL & vapour pressure $\times$ aqueous solubility & \textasciitilde\num{36000} & paired: $\Kwa = C_{\mathrm{aq}}^{\mathrm{sat}}/C_{\mathrm{gas}}^{\mathrm{sat}}$ \\
DGT, DGTX, DGB, DGV & transfer / vaporisation free energies & --- & excluded (not hydration) \\
PKA, KA, LAQACTCO & p$K_a$, association, activity coeff. & --- & excluded \\
\bottomrule
\end{tabular}
\caption{Measurement-type codes and their conversion routes. Henry's-law constants are
converted to the dimensionless air/water partition \Kwa\ accounting for the reported
convention (pressure/concentration, concentration/pressure, mole-fraction and dimensionless
forms); vapour pressures and solubilities are paired per molecule.}
\label{tab:process}
\end{table}

Harmonization yields \num{12027} observations over \num{2473} molecules
(\num{1518} with $\ge 2$ observations). We validated each route by its
per-observation residual against FreeSolv on the overlap, which exposed two systematic
conversion errors that we corrected: (i) the mole-fraction Henry route was biased by exactly
$-RT\ln 10^3 = +\SI{4.09}{\kcal}$, a dropped litre-to-cubic-metre factor between aqueous and
gas-phase concentrations; and (ii) the nominally ``water/gas'' dimensionless ratios are in
fact stored as the standard air/water Henry constant and required inversion (verified
empirically: the ratio of $1/\text{value}$ to the FreeSolv-implied \Kwa\ is $0.963$). After
correction, the per-molecule median of the harmonized observations agrees with FreeSolv to a
mean absolute error (MAE) of \SI{0.40}{\kcal} with Pearson $R=0.948$ (improved from
\SI{0.79}{\kcal}, $R=0.874$ before the fixes). The direct free-energy records
(FEOH, DGS) show essentially zero median residual and constitute the clean reference; the
Henry and pairing routes are the noisier, conflict-generating sources.

\section{Methods}

\subsection{Homoset: a proportional-similarity consensus gate}
The Homoset method treats each measurement type as a data \emph{source} and the direct
free-energy records as the trusted \emph{reference}. It proceeds in two stages.

\paragraph{Stage 1 --- source alignment.} For a non-reference source $s$, let $d_1,\dots,d_n$
be the per-molecule differences between $s$ and the reference over their overlap. Define an
offset and a scatter,
\begin{equation}
\delta_s = \operatorname{median}(d_i), \qquad
\mathrm{RMS}_s = \sqrt{\tfrac{1}{n}\textstyle\sum_i (d_i-\delta_s)^2},
\end{equation}
a noise/amplitude scale $L$, and the proportional-similarity index and its critical value
\begin{equation}
\mathrm{PS}_s = 1 - \frac{\mathrm{RMS}_s}{L}, \qquad
\mathrm{PS}_{\mathrm{crit}}(n,\alpha) = 1 - \sqrt{\frac{\chi^2_{n-1}(\alpha)}{3n}}.
\label{eq:ps}
\end{equation}
Source $s$ is \emph{accepted} into the consensus, and shifted by $-\delta_s$, iff
$\mathrm{PS}_s \ge \mathrm{PS}_{\mathrm{crit}}$; otherwise it is discarded. We use
$\alpha=0.01$. The noise scale is swept over two definitions: a \emph{data-adaptive} scale
$L = \max((\max d - \min d)/2,\, 0.5)$ (half the observed range of the differences), and a
\emph{fixed physical} scale $L=\SI{0.6}{\kcal}$ anchored to FreeSolv's typical experimental
uncertainty.

\paragraph{Stage 2 --- consensus and triage.} For each molecule the consensus is the median
over accepted, aligned observations; the molecule is flagged as \emph{conflicted} if that set
spans more than an outlier threshold (swept at $2.0$ and \SI{1.0}{\kcal}), yielding a direct
\texttt{use\_for\_training} decision.

\subsection{mixle: hierarchical partial pooling}
The mixle-style model is generative. With $y_{ij}$ the $j$-th observation of molecule $i$ from
source $s(ij)$,
\begin{equation}
y_{ij} \mid \mu_i \sim \mathcal N\!\big(\mu_i + b_{s(ij)},\, \sigma_{s(ij)}^2\big), \qquad
\mu_i \sim \mathcal N(\mu_0, \tau^2), \qquad b_{\mathrm{ref}} = 0 .
\label{eq:mixle}
\end{equation}
Anchoring the reference bias to zero fixes the absolute scale to the trusted free-energy
source. The per-source bias $b_s$, per-source noise $\sigma_s$, prior mean $\mu_0$, between-molecule
variance $\tau^2$, and the shrunken posteriors $\mu_i$ are estimated jointly by
expectation--maximization; the posterior mean of $\mu_i$ is a precision-weighted compromise
between its own (bias-corrected) observations and the global prior, so sparsely-measured
molecules are pooled toward $\mu_0$. A robust variant replaces the Gaussian observation model
with a Student-$t$ ($\nu=4$), realized as a scale mixture whose latent per-observation weights
down-weight outliers. Unlike the Homoset gate, mixle never rejects a source: it estimates the
source's systematic offset and \emph{propagates that correction to every molecule}, including
those with no reference measurement.

\section{Results}
We score each reconciled point estimate against FreeSolv on the 556-molecule overlap
(Table~\ref{tab:main}). Both methods beat the raw mean and median at all obs counts; the gap
widens on the 193 molecules the Homoset gate flags as conflicted, where reconciliation matters
most.

\begin{table}[t]
\centering
\small
\sisetup{table-format=1.3}
\begin{tabular}{@{}l S S S S S c@{}}
\toprule
estimator & {all MAE} & {all RMSE} & {$\ge$4-obs MAE} & {conflict MAE} & {conflict RMSE} & {bias} \\
\midrule
raw mean               & 0.570 & 1.226 & 0.406 & 1.134 & 1.698 & $+0.40$ \\
raw median             & 0.399 & 1.180 & 0.203 & 0.706 & 1.611 & $+0.23$ \\
Homoset, $L$ adaptive  & 0.392 & 1.159 & 0.202 & 0.697 & 1.592 & $+0.22$ \\
\textbf{Homoset, $L$ fixed 0.6} & \bfseries 0.283 & 0.979 & \bfseries 0.083 & 0.410 & 1.181 & $+0.15$ \\
\textbf{mixle, gaussian} & 0.308 & \bfseries 0.866 & 0.141 & \bfseries 0.382 & \bfseries 0.875 & \bfseries $+0.07$ \\
mixle, robust-$t$      & 0.276 & 0.949 & 0.116 & 0.388 & 1.152 & $+0.09$ \\
\bottomrule
\end{tabular}
\caption{Accuracy against FreeSolv on the overlap ($n=556$; \si{\kcal}). Bold marks the best
value in each column. The outlier threshold (2.0 vs.\ \SI{1.0}{\kcal}) affects only which
molecules are flagged, not the consensus value, so both settings score identically.}
\label{tab:main}
\end{table}

\paragraph{The noise scale governs the Homoset gate.} With the data-adaptive $L$, a few wildly
disagreeing molecules inflate the half-range to \SIrange{5.9}{10.3}{\kcal}, so
$\mathrm{PS}\approx 0.86 \gg \mathrm{PS}_{\mathrm{crit}}\approx 0.47$ and every noisy source is
admitted; the consensus barely improves on a plain median. With the fixed physical
$L=\SI{0.6}{\kcal}$, the scatter of the converted-source differences drives $\mathrm{PS}<0$, so
all converted sources are rejected and only the direct free-energy records survive. The gate
thus becomes a hard source-quality filter, and it is this strict setting that attains the best
median-based accuracy. The choice of $L$ is therefore the single most consequential parameter
of the method.

\paragraph{mixle corrects rather than rejects.} The Gaussian model estimates systematic source
biases of $+1.18$, $+1.25$ and \SI{+0.30}{\kcal} for the gas/water Henry, vapour-pressure/solubility
pairing, and water/gas Henry sources respectively (relative to the free-energy reference), and
subtracts them. This yields the lowest RMSE and the smallest bias of any method, and --- because
the bias is shared across molecules --- it is the only method that improves molecules lacking a
reference measurement. For example, a molecule with two Henry observations and no free-energy
record (InChIKey \texttt{CVXBEEMKQHEXEN}) has FreeSolv value \SI{-9.45}{\kcal}; the raw median and
the Homoset consensus are both stuck at \SI{-3.38}{\kcal}, whereas mixle-robust recovers
\SI{-9.22}{\kcal}. Conversely, for contaminated molecules that \emph{do} have a reference, the
strict Homoset gate and the robust model both recover the truth almost exactly (e.g.\
\texttt{ZWRUINPWMLAQRD}, 9 observations: truth \SI{-3.88}{\kcal}, raw median \SI{+1.63}{\kcal},
both methods \SI{-3.88}{\kcal}).

\section{Discussion}
The two methods embody a general trade-off in multi-source curation. The Homoset gate makes
discrete, auditable decisions --- accept or reject each source, flag or keep each molecule --- and
so produces exactly the artifact a curator wants: a defensible \texttt{use\_for\_training} label
per molecule. Its accuracy, however, is entirely contingent on choosing the noise scale $L$ to
match the true measurement error; with a data-adaptive $L$ it is scarcely better than a median.
The hierarchical model makes continuous, joint decisions --- soft per-source bias correction and
soft per-observation down-weighting --- and so extracts more signal, keeps all data, needs no
threshold, and generalizes source corrections to unmeasured cases; but it emits no explicit
curation flag and is less transparent. A practical pipeline uses both: reconcile values with the
hierarchical model, then triage with the Homoset gate to decide which molecules a human should
re-examine.

A caveat on evaluation: the FreeSolv overlap is enriched in molecules that possess a clean
free-energy record, which favours the strict Homoset setting (it can always fall back to that
record). On the roughly \num{1900} non-overlap molecules whose only data are Henry or pairing
measurements, the hierarchical model's ability to transfer a learned source bias should be more
valuable. Extending an independent validation to that regime is the natural next step, as is
modelling temperature and standard-state as explicit effects rather than pre-correcting them.

\section{Conclusion}
We converted the heterogeneous GuthrieSolv dump into \num{12027} standardized hydration
free-energy observations and used its redundancy to compare a proportional-similarity consensus
gate with a hierarchical partial-pooling model. Both recover the signal that naïve pooling
misses. The gate is the better curator under a physically-anchored noise scale; the hierarchical
model is the better reconciler and the more robust on hard, reference-poor molecules. The
harmonized dataset, conversion code, and both estimators are released to support continued
community curation of the database.

\section*{Data and code availability}
Harmonized observations, per-molecule estimates from all methods, full metrics, and the three
scripts (\texttt{harmonize\_guthrie.py}, \texttt{compare\_methods.py},
\texttt{validate\_vs\_freesolv.py}) are in \texttt{homoset\_vs\_mixle/} on the \texttt{homoset}
branch of the repository.

\begin{thebibliography}{9}
\bibitem{freesolv} D.~L.~Mobley and J.~P.~Guthrie, ``FreeSolv: a database of experimental and
calculated hydration free energies, with input files,'' \emph{J.\ Comput.-Aided Mol.\ Des.}
\textbf{28}, 711--720 (2014).
\bibitem{guthriesolv} MobleyLab, \emph{GuthrieSolv: The Guthrie Hydration Free Energy Database},
\url{https://github.com/MobleyLab/GuthrieSolv}.
\end{thebibliography}

\end{document}
